% page 179 of the facsimile (image p-178 of the scan)
% block 170.2 x 242.0 mm ; left margin 31.8 mm
\pgc{10.160mm}{0.000mm}{
\l{\cel{56}169.}
\l{}
\l{}
\l{\sou{flirtar}. (\sou{netrans.})\cel{2}(\sou{aludante personi di sexui interdiferanta)}.}
\l{\cel{3}Manovrar (per vorti o per konduto pleziva) por seduktar. - DEFR.}
\l{}
\l{\sou{flitro}. Lamelo de oro, de arjento, o de kupro orizita od arjenti-}
\l{\cel{3}zita, per qua onu plubriligas broduri, galoni, rubandi, e c. -D.}
\l{}
\l{\sou{flogar}. (\sou{trans.}) Batar per ledro-bendo o kordeto, agitante la}
\l{\cel{3}mancho ye qua fixigesas un ek la extremaji. - E.}
\l{}
\l{\sou{floko}. I. Mikra tufo lejera ek lano, silko, kotono. - II. Mikra}
\l{\cel{3}maso, poke densa.\cel{2}- DEFIR.}
\l{}
\l{\sou{floro}. (\sou{bot.}) Parto di la planto, qua su developas pos la folii,}
\l{\cel{3}desklozas sua korolo ofte odoroza, ornita ye kolori plu o min}
\l{\cel{3}briloza, e, pos existo provizora, remplasesas da la frukto.-}
\l{\cel{3}DEFIRS.}
\l{}
\l{\sou{florino}. Moneto-peco (olim ora, nun arjenta) uzata en Nederlando,}
\l{\cel{3}e c. - DEFIRS.}
\l{}
\l{\sou{florono}. I. Ornivo imitanta la formo di floro. - II. (\sou{arkitekt.})}
\l{\cel{3}Ornivo skultura qua reprezentas floro, folio. - FIS.}
\l{}
\l{\sou{floso}. Apendico membranoza per qua su movas en aquo la fishi e}
\l{\cel{3}kelka altra animali aquala. - D.}
\l{}
\l{\sou{floto}. Asemblajo plu o min granda de navi militala o komercala}
\l{\cel{3}qui navigas kune, ed havas un iro-skopo (generale komandata da}
\l{\cel{3}un chefo).\cel{3}EF.}
\l{}
\l{\sou{flotacar}. (\sou{netrans.}) Portesar sur liquido segun la arbitrio di la}
\l{\cel{3}ondi, di la fluo. - DEFIS.}
\l{}
\l{\sou{flotilio}. Floto de mikra navi. - DEFIRS.}
\l{}
\l{\sou{floxo}. (\sou{bot}.) Planto di qua la infloresenco esas piramidatra.}
\l{\cel{3}L.\cel{1}\sou{phlox}.}
\l{}
\l{\sou{flozelo}. Buro de silko, reziduo de la kokoni desvolvita. - EF.}
\l{}
\l{\sou{fluar}. (\sou{netrans}.)(\sou{aludante liquido}). Irar de loko ad altra loko,}
\l{\cel{3}per movo kontinua. - (\sou{anke metaf}.). - DEF.}
\l{}
\l{\sou{flugar}.\cel{2}(\sou{netrans}.) Sustenigar su e movar su en aero atmosfero}
\l{\cel{3}per ali (naturala od artificala). (ex.: la uceli e la maxim}
\l{\cel{3}multa insekti ; la aviacili : aeroplani, helikopteri). - DE.}
\l{}
\l{\sou{fluida}. Di qua la molekuli, poke inter-adheranta, ne rezistas}
\l{\cel{3}mem nur mikra esforco qua tendencas a diplasar li. - DEFIRS.}
\l{}
\l{fluktuar. (\sou{netrans.}) Quaze ondifar en aero atmosfero, segun la}
\l{\cel{3}arbitrio di vento o di altra impulso variiva. (kom ex.:}
\l{\cel{3}standardo, hari.) - (\sou{anke pri objekto imersita}). - DEFIRS.}
}
